\subsection*{Tileflow 程序：编排数据搬运}

下面把迭代、分块选择和 DMA 放在同一个 Tileflow 程序中，观察数据从输入到输出的过程。

完整代码如下：

\begin{lstlisting}[language=croqtile,caption={tileflow 程序}]
__co__ s32 [6, 17, 128] ele_add(s32 [6, 17, 128] lhs, s32 [6, 17, 128] rhs) {
  s32 [lhs.span] output; // Use same shape as lhs
  parallel p by 6 {
    with index in [17, 4] { // Tiling factors
      foreach index {
        lhs_load = dma.copy lhs.chunkat(p, index) => local;
        rhs_load = dma.copy rhs.chunkat(p, index) => local;
        local s32 [lhs_load.span] l1_out;
        // Call kernel with loaded data
        call kernel(lhs_load.data, rhs_load.data, l1_out, |lhs_load.span|);
        // Store result back to output
        dma.copy l1_out => output.chunkat(p, index);
      }
    }
  }
  return output;
}
\end{lstlisting}

\co@__co__@ 函数接收形状为 \co@[6, 17, 128]@ 的两个 \co@s32@ 输入，并声明同形状的输出。

\co@parallel p by 6@ 声明 6 个独立实例，每个实例处理输入的一个外层区域。

\co@with-in@ 将 \co@index@ 绑定到范围 \co@[17, 4]@。它是二维有界整数元组，供内部循环使用。

\co@foreach index@ 对应下面的嵌套循环：

\begin{lstlisting}[language=croqtile,caption={等价的 C++ 循环}]
for (int x = 0; x < 17; x++)  // assume 'x' represents the 1st element of 'index'
  for (int y = 0; y < 4; y++) { ... }  // and 'y' represents the 2nd element of 'index'
\end{lstlisting}

循环中的 \co@dma.copy@ 负责搬运当前分块：

\begin{itemize}[leftmargin=1.6em]
\item \co@lhs_load@：传输结果句柄，用于引用目的缓冲区。
\item \co@dma.copy@：拷贝数据，\co@=>@ 左边是源，右边是目的。
\item \co@local@：让编译器分配线程私有的目的缓冲区。
\item \co@lhs.chunkat(p, index)@：根据范围 \co@6@、\co@17@、\co@4@ 选取 \co@[1, 1, 32]@ 分块，具体位置由索引当前值决定。例如索引 \co@(1, 16, 2)@ 选中对应的分块。
\end{itemize}

分块方式如下图所示。

\begin{figure}[htbp]\centering
\includegraphics[width=\linewidth,height=0.56\textheight,keepaspectratio]{f14-chunkat_3d_light.png}
\caption{形状 [6, 17, 128] 按范围 [6, 17, 4] 划分后，每块为 [1, 1, 32]。分块 (1, 16, 2) 从元素 (1, 16, 64) 开始。}
\end{figure}

\co@lhs@ 和 \co@rhs@ 的对应分块分别搬入设备的 \co@local@ 存储，供当前实例计算。

\co@local s32 [lhs_load.span] l1_out;@ 声明同分块形状的输出缓冲区。\co@call kernel@ 在其中完成计算，随后 DMA 将结果搬回主机侧输出区域。

每个实例迭代 17 × 4 次，每次处理一个 \co@[1, 1, 32]@ 分块。全部实例完成后，函数把输出交给主机调用方。
